% fig2_algorithm_flowchart.tex —— MPC + 走廊 + softmin 框架流程图
% 遵守 §9 红线：走廊代价仅「垂直偏离 hinge」一项（半宽 = 拍半径 0.12 m），
% 速度方向/法向量对齐由时间窗 softmin 终端承担，不写入走廊注释。
% 用法：正文 % fig2_algorithm_flowchart.tex —— MPC + 走廊 + softmin 框架流程图
% 遵守 §9 红线：走廊代价仅「垂直偏离 hinge」一项（半宽 = 拍半径 0.12 m），
% 速度方向/法向量对齐由时间窗 softmin 终端承担，不写入走廊注释。
% 用法：正文 % fig2_algorithm_flowchart.tex —— MPC + 走廊 + softmin 框架流程图
% 遵守 §9 红线：走廊代价仅「垂直偏离 hinge」一项（半宽 = 拍半径 0.12 m），
% 速度方向/法向量对齐由时间窗 softmin 终端承担，不写入走廊注释。
% 用法：正文 % fig2_algorithm_flowchart.tex —— MPC + 走廊 + softmin 框架流程图
% 遵守 §9 红线：走廊代价仅「垂直偏离 hinge」一项（半宽 = 拍半径 0.12 m），
% 速度方向/法向量对齐由时间窗 softmin 终端承担，不写入走廊注释。
% 用法：正文 \input{figures/fig2_algorithm_flowchart.tex}
% 配色与结果图统一（chat9 审稿意见）：走廊蓝 = 结果图 corridor-only 蓝 (0,114,178)
% —— 旧版 blue!18 偏紫；softmin 琥珀 = softmin-only 琥珀 (230,159,0)
\definecolor{corridorblue}{RGB}{0,114,178}
\definecolor{softminamber}{RGB}{230,159,0}
\begin{figure}[t]
\centering
\begin{tikzpicture}[node distance=0.38cm, >=stealth,
    block/.style={rectangle, draw, fill=white, rounded corners=1pt,
                  text width=4.8cm, minimum height=0.62cm, font=\footnotesize,
                  align=center, inner sep=2pt},
    corridorcost/.style={rectangle, draw=corridorblue, fill=corridorblue!18,
                 rounded corners=1pt,
                 text width=4.8cm, minimum height=0.62cm, font=\footnotesize,
                 align=center, inner sep=2pt},
    softmincost/.style={rectangle, draw=softminamber, fill=softminamber!22,
                 rounded corners=1pt,
                 text width=4.8cm, minimum height=0.62cm, font=\footnotesize,
                 align=center, inner sep=2pt}]

% chat9 审稿意见「精简 Fig. 2」：观测与预测合并为一个块，节点间距收窄
% chat12 审稿意见：两 cost 块的归属词（running / terminal cost）提到行首成为
% 最醒目信息，参数细节（r_racket、β=5）删除（正文/图注已有）
\node[block] (obs) {Observe ball state, predict trajectory\\and candidate windows};
\node[corridorcost, below=of obs] (corr) {Running cost: corridor hinge term};
\node[softmincost, below=of corr] (soft) {Terminal cost: candidate-window softmin};
\node[block, below=of soft] (ilqr) {iLQR backward $+$ forward pass (MPC)};
\node[block, below=of ilqr] (exec) {Execute with safety filter\\(joint/TCP limits, midline plane)};

\draw[->] (obs) -- (corr);
\draw[->] (corr) -- (soft);
\draw[->] (soft) -- (ilqr);
\draw[->] (ilqr) -- (exec);
\draw[->] (exec.east) -- ++(0.85,0) |- (obs.east)
    node[pos=0.35, right, font=\footnotesize, align=left] {replan\\interval};

\end{tikzpicture}
\caption{MPC-iLQR framework. Shaded blocks are the two cost modifications; the
point-target baseline (\texttt{none}) drops both.}
\label{fig:algorithm_flowchart}
\end{figure}

% 配色与结果图统一（chat9 审稿意见）：走廊蓝 = 结果图 corridor-only 蓝 (0,114,178)
% —— 旧版 blue!18 偏紫；softmin 琥珀 = softmin-only 琥珀 (230,159,0)
\definecolor{corridorblue}{RGB}{0,114,178}
\definecolor{softminamber}{RGB}{230,159,0}
\begin{figure}[t]
\centering
\begin{tikzpicture}[node distance=0.38cm, >=stealth,
    block/.style={rectangle, draw, fill=white, rounded corners=1pt,
                  text width=4.8cm, minimum height=0.62cm, font=\footnotesize,
                  align=center, inner sep=2pt},
    corridorcost/.style={rectangle, draw=corridorblue, fill=corridorblue!18,
                 rounded corners=1pt,
                 text width=4.8cm, minimum height=0.62cm, font=\footnotesize,
                 align=center, inner sep=2pt},
    softmincost/.style={rectangle, draw=softminamber, fill=softminamber!22,
                 rounded corners=1pt,
                 text width=4.8cm, minimum height=0.62cm, font=\footnotesize,
                 align=center, inner sep=2pt}]

% chat9 审稿意见「精简 Fig. 2」：观测与预测合并为一个块，节点间距收窄
% chat12 审稿意见：两 cost 块的归属词（running / terminal cost）提到行首成为
% 最醒目信息，参数细节（r_racket、β=5）删除（正文/图注已有）
\node[block] (obs) {Observe ball state, predict trajectory\\and candidate windows};
\node[corridorcost, below=of obs] (corr) {Running cost: corridor hinge term};
\node[softmincost, below=of corr] (soft) {Terminal cost: candidate-window softmin};
\node[block, below=of soft] (ilqr) {iLQR backward $+$ forward pass (MPC)};
\node[block, below=of ilqr] (exec) {Execute with safety filter\\(joint/TCP limits, midline plane)};

\draw[->] (obs) -- (corr);
\draw[->] (corr) -- (soft);
\draw[->] (soft) -- (ilqr);
\draw[->] (ilqr) -- (exec);
\draw[->] (exec.east) -- ++(0.85,0) |- (obs.east)
    node[pos=0.35, right, font=\footnotesize, align=left] {replan\\interval};

\end{tikzpicture}
\caption{MPC-iLQR framework. Shaded blocks are the two cost modifications; the
point-target baseline (\texttt{none}) drops both.}
\label{fig:algorithm_flowchart}
\end{figure}

% 配色与结果图统一（chat9 审稿意见）：走廊蓝 = 结果图 corridor-only 蓝 (0,114,178)
% —— 旧版 blue!18 偏紫；softmin 琥珀 = softmin-only 琥珀 (230,159,0)
\definecolor{corridorblue}{RGB}{0,114,178}
\definecolor{softminamber}{RGB}{230,159,0}
\begin{figure}[t]
\centering
\begin{tikzpicture}[node distance=0.38cm, >=stealth,
    block/.style={rectangle, draw, fill=white, rounded corners=1pt,
                  text width=4.8cm, minimum height=0.62cm, font=\footnotesize,
                  align=center, inner sep=2pt},
    corridorcost/.style={rectangle, draw=corridorblue, fill=corridorblue!18,
                 rounded corners=1pt,
                 text width=4.8cm, minimum height=0.62cm, font=\footnotesize,
                 align=center, inner sep=2pt},
    softmincost/.style={rectangle, draw=softminamber, fill=softminamber!22,
                 rounded corners=1pt,
                 text width=4.8cm, minimum height=0.62cm, font=\footnotesize,
                 align=center, inner sep=2pt}]

% chat9 审稿意见「精简 Fig. 2」：观测与预测合并为一个块，节点间距收窄
% chat12 审稿意见：两 cost 块的归属词（running / terminal cost）提到行首成为
% 最醒目信息，参数细节（r_racket、β=5）删除（正文/图注已有）
\node[block] (obs) {Observe ball state, predict trajectory\\and candidate windows};
\node[corridorcost, below=of obs] (corr) {Running cost: corridor hinge term};
\node[softmincost, below=of corr] (soft) {Terminal cost: candidate-window softmin};
\node[block, below=of soft] (ilqr) {iLQR backward $+$ forward pass (MPC)};
\node[block, below=of ilqr] (exec) {Execute with safety filter\\(joint/TCP limits, midline plane)};

\draw[->] (obs) -- (corr);
\draw[->] (corr) -- (soft);
\draw[->] (soft) -- (ilqr);
\draw[->] (ilqr) -- (exec);
\draw[->] (exec.east) -- ++(0.85,0) |- (obs.east)
    node[pos=0.35, right, font=\footnotesize, align=left] {replan\\interval};

\end{tikzpicture}
\caption{MPC-iLQR framework. Shaded blocks are the two cost modifications; the
point-target baseline (\texttt{none}) drops both.}
\label{fig:algorithm_flowchart}
\end{figure}

% 配色与结果图统一（chat9 审稿意见）：走廊蓝 = 结果图 corridor-only 蓝 (0,114,178)
% —— 旧版 blue!18 偏紫；softmin 琥珀 = softmin-only 琥珀 (230,159,0)
\definecolor{corridorblue}{RGB}{0,114,178}
\definecolor{softminamber}{RGB}{230,159,0}
\begin{figure}[t]
\centering
\begin{tikzpicture}[node distance=0.38cm, >=stealth,
    block/.style={rectangle, draw, fill=white, rounded corners=1pt,
                  text width=4.8cm, minimum height=0.62cm, font=\footnotesize,
                  align=center, inner sep=2pt},
    corridorcost/.style={rectangle, draw=corridorblue, fill=corridorblue!18,
                 rounded corners=1pt,
                 text width=4.8cm, minimum height=0.62cm, font=\footnotesize,
                 align=center, inner sep=2pt},
    softmincost/.style={rectangle, draw=softminamber, fill=softminamber!22,
                 rounded corners=1pt,
                 text width=4.8cm, minimum height=0.62cm, font=\footnotesize,
                 align=center, inner sep=2pt}]

% chat9 审稿意见「精简 Fig. 2」：观测与预测合并为一个块，节点间距收窄
% chat12 审稿意见：两 cost 块的归属词（running / terminal cost）提到行首成为
% 最醒目信息，参数细节（r_racket、β=5）删除（正文/图注已有）
\node[block] (obs) {Observe ball state, predict trajectory\\and candidate windows};
\node[corridorcost, below=of obs] (corr) {Running cost: corridor hinge term};
\node[softmincost, below=of corr] (soft) {Terminal cost: candidate-window softmin};
\node[block, below=of soft] (ilqr) {iLQR backward $+$ forward pass (MPC)};
\node[block, below=of ilqr] (exec) {Execute with safety filter\\(joint/TCP limits, midline plane)};

\draw[->] (obs) -- (corr);
\draw[->] (corr) -- (soft);
\draw[->] (soft) -- (ilqr);
\draw[->] (ilqr) -- (exec);
\draw[->] (exec.east) -- ++(0.85,0) |- (obs.east)
    node[pos=0.35, right, font=\footnotesize, align=left] {replan\\interval};

\end{tikzpicture}
\caption{MPC-iLQR framework. Shaded blocks are the two cost modifications; the
point-target baseline (\texttt{none}) drops both.}
\label{fig:algorithm_flowchart}
\end{figure}
